% !TeX program = lualatex
% ============================================================
% appendixA.tex — appA 唯一內容源（2026-08-09 起；LaTeX 統一 U1）
% 沿革：升格自 dist/appA/appendixA.tex（P1 首轉：make_dist.py 自 fragment 轉換；
%   fragment 樹已凍結，改課文一律改本檔。拍板見 ../../KICKOFF-latex-unification.md）
%   樣式源＝../../template/calcbook.sty（M-B1 拍板模板）
% 編譯（本資料夾為 CWD；aux 進 build/、成品 PDF 移入 ../../dist/appA/）：
%   latexmk -lualatex -auxdir=../../build/aux-appA appendixA.tex
% ============================================================
\documentclass[a4paper,12pt,oneside]{memoir}
\makeatletter\def\input@path{{../../template/}}\makeatother
\def\cbfontsdir{../../template/fonts/inter/}
\usepackage{calcbook}
% 圖：emitter 射 `<ch>/<stem>`（見 convert.py），故 graphicspath 指向 chapters/ 上層。
% 模板內的 {../figs/} 是 chapters/<ch>/ 為 CWD 時的路徑，dist/<ch>/ 需這一行覆寫。
\graphicspath{{../../chapters/}}
\begin{document}
\appendixopener{Appendix A}{Precalculus Toolbox}
\begin{lead}
This appendix collects the precalculus tools you will want on hand but that many precalculus courses now treat briefly or omit: the reciprocal trigonometric functions (§A.1), the product-to-sum and sum-to-product identities (§A.2), sigma notation, power sums, and the geometric series (§A.3), partial fractions (§A.4), a quick tour of the conic sections (§A.5), and absolute value and inequalities (§A.6). Some are pointed to from the main text at the moment they are first needed; others are here as a handy reference to reach for when you want them. The sections are independent of one another; read each when the main text sends you, or straight through as a warm-up. Nothing here uses calculus, with one small exception flagged in §A.3, where an infinite sum is read as a limit.
\end{lead}
\parahead{By the end of this appendix, you will be able to:}
\begin{objectives}
  \item evaluate the reciprocal trigonometric functions \(\sec\), \(\csc\), and \(\cot\), and apply their Pythagorean identities with the correct signs;
  \item derive every product-to-sum and sum-to-product identity from the angle-sum formulas — and recover any one you forget;
  \item manipulate sums in sigma notation, evaluate the power sums \(\sum k\), \(\sum k^{2}\), and \(\sum k^{3}\), and sum a geometric series, finite or infinite;
  \item split a proper rational function into partial fractions;
  \item read off each conic section’s key features — centre, vertices, foci, and axes — from its standard equation;
  \item unpack an absolute-value inequality \(|x - a| < \delta\) into an interval, and apply the triangle inequality.
\end{objectives}
\sechead{A.1}{The Reciprocal Trigonometric Functions}
The three primary trigonometric functions \(\sin\), \(\cos\), and \(\tan\) have three companions — two honest reciprocals and one inverted ratio. They carry no new information — every fact about them is a fact about \(\sin\) and \(\cos\) in disguise. Yet calculus speaks them fluently: the derivative of \(\tan x\) turns out to be \(\sec^{2} x\), the inverse secant needs a principal range of its own, and integral tables are full of secants and cosecants. This section gives them an honest introduction.

\begin{envdefinition}{Definition}{A.1}{}
For any angle \(\theta\) at which the denominator is nonzero, the \emph{secant}, \emph{cosecant}, and \emph{cotangent} functions are defined by \[ \sec\theta = \frac{1}{\cos\theta}, \qquad \csc\theta = \frac{1}{\sin\theta}, \qquad \cot\theta = \frac{\cos\theta}{\sin\theta}. \]

\begin{informal}
Informally: secant is the reciprocal of cosine, cosecant is the reciprocal of sine, and cotangent flips the ratio that defines tangent. The names cross over historically — \emph{se}cant with \emph{co}sine, \emph{co}secant with sine — and the pairing simply has to be memorized.
\end{informal}
\end{envdefinition}
In a right triangle, the definitions invert the familiar ratios. Since \(\cos\theta\) is adjacent over hypotenuse, \(\sec\theta\) is hypotenuse over adjacent; since \(\sin\theta\) is opposite over hypotenuse, \(\csc\theta\) is hypotenuse over opposite; and \(\cot\theta\) is adjacent over opposite. Special-angle values follow instantly from the primary ones: for instance \(\sec\tfrac{\pi}{3} = 1/\cos\tfrac{\pi}{3} = 2\), and \(\cot\tfrac{\pi}{4} = 1\). Each function is undefined exactly where its denominator vanishes — \(\sec\) and \(\tan\) at odd multiples of \(\tfrac{\pi}{2}\), \(\csc\) and \(\cot\) at multiples of \(\pi\) — and near those angles its values blow up, producing the vertical asymptotes on their graphs.

\begin{envcaution}{Caution}{}{}
Cotangent is defined as \(\cos\theta/\sin\theta\), \emph{not} as \(1/\tan\theta\). On most of the domain the two agree, but not everywhere: at \(\theta = \tfrac{\pi}{2}\) we have \(\cot\tfrac{\pi}{2} = \cos\tfrac{\pi}{2}/\sin\tfrac{\pi}{2} = 0/1 = 0\), a perfectly good value, while \(1/\tan\tfrac{\pi}{2}\) is meaningless — \(\tan\tfrac{\pi}{2}\) is undefined. The quotient \(\cos/\sin\) is the robust definition; \(1/\tan\) is a shortcut valid only where \(\tan\theta\) is defined and nonzero.
\end{envcaution}
Dividing the Pythagorean identity \(\sin^{2}\theta + \cos^{2}\theta = 1\) by \(\cos^{2}\theta\) or by \(\sin^{2}\theta\) produces the two companion identities that calculus uses constantly.

\begin{envproposition}{Proposition}{A.1}{Pythagorean companions}
\[ \begin{aligned} 1 + \tan^{2}\theta &= \sec^{2}\theta \qquad (\cos\theta \neq 0), \\ 1 + \cot^{2}\theta &= \csc^{2}\theta \qquad (\sin\theta \neq 0). \end{aligned} \]
\end{envproposition}
\begin{envproof}{Proof}{}{}
Divide \(\sin^{2}\theta + \cos^{2}\theta = 1\) through by \(\cos^{2}\theta\): \[ \frac{\sin^{2}\theta}{\cos^{2}\theta} + 1 = \frac{1}{\cos^{2}\theta}, \qquad\text{that is,}\qquad \tan^{2}\theta + 1 = \sec^{2}\theta. \] Dividing instead by \(\sin^{2}\theta\) gives \(1 + \cot^{2}\theta = \csc^{2}\theta\). \qedmark
\end{envproof}
\begin{workedexample}
\begin{envexample}{Example}{A.1}{}
Evaluate the following.

\begin{enumerate}
  \item \(\sec\!\left(\tfrac{2\pi}{3}\right)\) and \(\cot\!\left(-\tfrac{\pi}{4}\right)\).
  \item Given that \(\tan\theta = 3\) and \(\theta\) lies in the third quadrant, find \(\sec\theta\) and \(\csc\theta\).
\end{enumerate}
\end{envexample}
\begin{envsolution}{Solution}{}{}
\begin{enumerate}
  \item Since \(\cos\tfrac{2\pi}{3} = -\tfrac{1}{2}\), we get \(\sec\tfrac{2\pi}{3} = 1/(-\tfrac{1}{2}) = -2\). Since \(\cos(-\tfrac{\pi}{4}) = \tfrac{\sqrt{2}}{2}\) and \(\sin(-\tfrac{\pi}{4}) = -\tfrac{\sqrt{2}}{2}\), we get \(\cot(-\tfrac{\pi}{4}) = \frac{\sqrt{2}/2}{-\sqrt{2}/2} = -1\).
  \item By Proposition A.1, \(\sec^{2}\theta = 1 + \tan^{2}\theta = 1 + 9 = 10\), so \(\sec\theta = \pm\sqrt{10}\). In the third quadrant \(\cos\theta < 0\), hence \(\sec\theta < 0\), and therefore \(\sec\theta = -\sqrt{10}\). For the cosecant, \(\cot\theta = 1/\tan\theta = \tfrac{1}{3}\) (legitimate here, since \(\tan\theta = 3\) is defined and nonzero), so \(\csc^{2}\theta = 1 + \cot^{2}\theta = 1 + \tfrac{1}{9} = \tfrac{10}{9}\) and \(\csc\theta = \pm\tfrac{\sqrt{10}}{3}\). In the third quadrant \(\sin\theta < 0\), so \(\csc\theta = -\tfrac{\sqrt{10}}{3}\). The squaring in Proposition A.1 forgets all signs; the quadrant restores them.
\end{enumerate}
\end{envsolution}
\end{workedexample}
That is the whole toolkit: three definitions, the right-triangle readings, and the two Pythagorean companions — every fact about \(\sec\), \(\csc\), and \(\cot\) reduced to a fact about \(\sin\) and \(\cos\).


\sechead{A.2}{The Product-to-Sum and Sum-to-Product Identities}
One algebraic step in this book carries more weight than any other: rewriting the difference \(\sin(x+h) - \sin x\) as a \emph{product}, which is what makes the derivative of the sine function computable. The identity that does it belongs to a family — the product-to-sum and sum-to-product identities — that older curricula drilled and most current precalculus courses omit. The good news is that the whole family is nothing but the angle-sum formulas, added and subtracted in pairs. If you learn one thing from this section, learn how to \emph{rebuild} the identities; the memorizing then becomes optional.

Everything starts from the four angle-sum and angle-difference formulas, which we assume and record here for reference: \[ \begin{aligned} \sin(u + v) &= \sin u \cos v + \cos u \sin v, \\ \sin(u - v) &= \sin u \cos v - \cos u \sin v, \\ \cos(u + v) &= \cos u \cos v - \sin u \sin v, \\ \cos(u - v) &= \cos u \cos v + \sin u \sin v. \end{aligned} \]

Add or subtract these in matching pairs and half of each right-hand side cancels, leaving a single product.

\begin{envproposition}{Proposition}{A.2}{Product-to-sum identities}
\[ \begin{aligned} \sin u \cos v &= \tfrac{1}{2}\bigl[\sin(u + v) + \sin(u - v)\bigr], \\ \cos u \cos v &= \tfrac{1}{2}\bigl[\cos(u + v) + \cos(u - v)\bigr], \\ \sin u \sin v &= \tfrac{1}{2}\bigl[\cos(u - v) - \cos(u + v)\bigr]. \end{aligned} \]
\end{envproposition}
\begin{envproof}{Proof}{}{}
Add the two sine formulas: the \(\cos u \sin v\) terms cancel, leaving \(\sin(u+v) + \sin(u-v) = 2 \sin u \cos v\); divide by \(2\). Add the two cosine formulas: the \(\sin u \sin v\) terms cancel, giving \(\cos(u+v) + \cos(u-v) = 2 \cos u \cos v\). Subtract the cosine formulas instead — \(\cos(u-v) - \cos(u+v)\) — and the \(\cos u \cos v\) terms cancel, giving \(2 \sin u \sin v\). \qedmark
\end{envproof}
Setting \(u = v = \theta\) in these product-to-sum identities pays an immediate dividend. The middle line becomes \(\cos^{2}\theta = \tfrac{1}{2}\bigl[\cos 2\theta + \cos 0\bigr]\) and the last becomes \(\sin^{2}\theta = \tfrac{1}{2}\bigl[\cos 0 - \cos 2\theta\bigr]\); since \(\cos 0 = 1\), these are the \emph{power-reduction} formulas \[ \cos^{2}\theta = \frac{1 + \cos 2\theta}{2}, \qquad \sin^{2}\theta = \frac{1 - \cos 2\theta}{2}, \] which trade a \emph{squared} trigonometric function for a plain one at twice the angle. (The first line, \(\sin\theta \cos\theta = \tfrac{1}{2}\sin 2\theta\), is the double-angle formula for sine read the same way.) These are the forms integration reaches for to handle \(\sin^{2}\) and \(\cos^{2}\) when the techniques of integration arrive.

The sum-to-product identities are the same cancellations read in the other direction. The only new idea is a change of names: given a sum like \(\sin A + \sin B\), arrange for the two angles to read \(u + v\) and \(u - v\) by choosing \(u\) and \(v\) so that \(u + v = A\) and \(u - v = B\). Adding these two equations gives \(2u = A + B\) and subtracting them gives \(2v = A - B\), so \(u\) is the \emph{average} and \(v\) the \emph{half-difference}: \[ u = \frac{A + B}{2}, \qquad v = \frac{A - B}{2}. \]

\begin{envproposition}{Proposition}{A.3}{Sum-to-product identities}
\[ \begin{aligned} \sin A + \sin B &= 2 \sin\frac{A + B}{2} \cos\frac{A - B}{2}, \\ \sin A - \sin B &= 2 \cos\frac{A + B}{2} \sin\frac{A - B}{2}, \\ \cos A + \cos B &= 2 \cos\frac{A + B}{2} \cos\frac{A - B}{2}, \\ \cos A - \cos B &= -2 \sin\frac{A + B}{2} \sin\frac{A - B}{2}. \end{aligned} \]
\end{envproposition}
\begin{envproof}{Proof}{}{}
With \(u = \tfrac{A+B}{2}\) and \(v = \tfrac{A-B}{2}\) we have \(u + v = A\) and \(u - v = B\). Then, reading the cancellations in the previous proof from right to left, \[ \begin{aligned} \sin A + \sin B &= \sin(u + v) + \sin(u - v) \\ &= 2 \sin u \cos v = 2 \sin\frac{A + B}{2} \cos\frac{A - B}{2}, \end{aligned} \] and the other three lines follow the same way: \(\sin A - \sin B = \sin(u+v) - \sin(u-v) = 2 \cos u \sin v\) (this time subtracting the two sine expansions); \(\cos A + \cos B = 2 \cos u \cos v\); and \(\cos A - \cos B = \cos(u+v) - \cos(u-v) = -2 \sin u \sin v\). Note the minus sign in the last line — the only one in the family, and the easiest thing to misremember; the subtraction order \(\cos(u+v) - \cos(u-v)\) is where it comes from. \qedmark
\end{envproof}
\begin{envstrategy}{Strategy}{A.1}{Recovering an identity you forgot}
\begin{steps}
  \item Write down the two angle-sum (or angle-difference) expansions that contain the pieces you have, and add or subtract them so that the unwanted terms cancel.
  \item If your expression is a sum of two trig values at angles \(A\) and \(B\), substitute the average \(u = \tfrac{A+B}{2}\) and half-difference \(v = \tfrac{A-B}{2}\), so that \(A = u + v\) and \(B = u - v\).
\end{steps}
\begin{informal}
Rebuilding takes perhaps a minute; a memorized list of eight is easy to get wrong by a sign. Verifying is even cheaper — expand the product side with the angle-sum formulas and watch the cross terms cancel.
\end{informal}
\end{envstrategy}
\begin{workedexample}
\begin{envexample}{Example}{A.2}{}
Rewrite each expression as requested.

\begin{enumerate}
  \item Write \(\sin(x + h) - \sin x\) as a product.
  \item Write \(\sin 5x \, \sin 3x\) as a sum.
\end{enumerate}
\end{envexample}
\begin{envsolution}{Solution}{}{}
\begin{enumerate}
  \item Apply the sum-to-product identity for \(\sin A - \sin B\) with \(A = x + h\) and \(B = x\). The average is \(\tfrac{A+B}{2} = x + \tfrac{h}{2}\) and the half-difference is \(\tfrac{A-B}{2} = \tfrac{h}{2}\), so \[ \sin(x + h) - \sin x = 2 \cos\Bigl(x + \tfrac{h}{2}\Bigr) \sin\frac{h}{2}. \] This is precisely the rewriting the derivative of sine turns on: the difference of two sines becomes a product with the small quantity \(h/2\) isolated inside a single sine.
  \item Apply the third product-to-sum identity with \(u = 5x\) and \(v = 3x\): \[ \begin{aligned} \sin 5x \, \sin 3x &= \tfrac{1}{2}\bigl[\cos(5x - 3x) - \cos(5x + 3x)\bigr] \\ &= \tfrac{1}{2}\bigl[\cos 2x - \cos 8x\bigr]. \end{aligned} \] Turning products into sums is the direction integration will want: a sum of cosines is far easier to integrate than a product.
\end{enumerate}
\end{envsolution}
\end{workedexample}
Of the whole family, the identity for \(\sin A - \sin B\) in Proposition A.3 is the one worth keeping at your fingertips; the rest follow from the same two moves the moment you need them.


\sechead{A.3}{Sigma Notation, Power Sums, and the Geometric Series}
Before the geometric series itself, a word on the notation that names it and a few moves for handling sums cleanly. The capital-sigma symbol \(\sum\) abbreviates a sum: \(\sum_{k=1}^{n} a_{k}\) means \(a_{1} + a_{2} + \cdots + a_{n}\), the terms \(a_{k}\) added as the \emph{index} \(k\) climbs through the whole numbers from the lower limit to the upper. The index is only a placeholder — \(\sum_{k=1}^{n} a_{k}\) and \(\sum_{j=1}^{n} a_{j}\) denote the very same sum — and an upper limit of \(\infty\), as in \(\sum_{k=0}^{\infty} a_{k}\), signals a sum that never stops, to be read as the limit of its running totals in the sense settled just below.

Three moves are worth having at hand. \emph{Re-indexing} renames the counter without disturbing the sum: setting \(j = k - 1\) rewrites \(\sum_{k=1}^{n} a_{k}\) as \(\sum_{j=0}^{n-1} a_{j+1}\), the same terms relabelled. \emph{Peeling} detaches the first or last term, as in \(\sum_{k=0}^{n} a_{k} = a_{0} + \sum_{k=1}^{n} a_{k}\). And \emph{splitting at a cutoff} \(n_{0}\) breaks a sum into a leading part and a \emph{tail}, \[ \sum_{k=0}^{n} a_{k} \;=\; \sum_{k=0}^{n_{0}} a_{k} \;+\; \sum_{k=n_{0}+1}^{n} a_{k}, \] a standard way to control a series: bound its tail while the leading partial sum is set aside.

One special sum collapses outright. When each term is the \emph{difference} of consecutive values of a sequence \(b_{k}\), the interior pieces cancel in pairs: \[ \sum_{k=0}^{m} \bigl(b_{k} - b_{k+1}\bigr) = b_{0} - b_{m+1}, \] since every \(b_{k}\) with \(1 \le k \le m\) is added once and subtracted once. A sum that collapses this way is said to \emph{telescope}, after the way a spyglass folds shut — and it is the exact mechanism behind both the geometric-sum formula and the power sums below.

Those moves are all finite bookkeeping. Adding up \emph{infinitely} many terms is a genuinely new act, and most of its theory belongs to a later chapter of calculus. But one infinite sum is so simple, and so useful so early, that it deserves to be settled now: the geometric series, in which each term is a fixed multiple \(r\) of the one before. It earns its keep everywhere — the \emph{tail} of a series (its leftover terms past a cutoff) can often be trapped under a geometric tail and squeezed to nothing. Everything rests on one finite identity plus one limit.

\begin{envproposition}{Proposition}{A.4}{Finite geometric sum}
For any number \(r \neq 1\) and any positive integer \(m\), \[ 1 + r + r^{2} + \cdots + r^{m} = \frac{1 - r^{m+1}}{1 - r}. \]
\end{envproposition}
\begin{envproof}{Proof}{}{}
Write \(S = 1 + r + r^{2} + \cdots + r^{m}\) and multiply by \(1 - r\): \[ (1 - r)S = \bigl(1 + r + \cdots + r^{m}\bigr) - \bigl(r + r^{2} + \cdots + r^{m+1}\bigr) = 1 - r^{m+1}, \] because each power \(r, r^{2}, \ldots, r^{m}\) appears once in each bracket and cancels, leaving only the leading \(1\) and the final \(-r^{m+1}\) — the sum telescopes. Dividing by \(1 - r\) (legitimate, since \(r \neq 1\)) gives the identity. \qedmark
\end{envproof}
Now let the sum run on. If \(\lvert r \rvert < 1\), each multiplication by \(r\) shrinks the next term by a fixed factor, and the power \(r^{m+1}\) in the formula dies away: its size \(\lvert r \rvert^{m+1}\) drops below any positive threshold once \(m\) is large enough. In the arrow notation for a limit (here the whole-number index \(m\) runs off to infinity, rather than a variable approaching a finite point), \(r^{m+1} \to 0\) as \(m \to \infty\), so the running totals \[ S_{m} = \frac{1 - r^{m+1}}{1 - r} \qquad\text{settle toward}\qquad \frac{1 - 0}{1 - r} = \frac{1}{1 - r}. \] The value of an infinite sum \emph{means} exactly this: the limit of its running totals. (The running totals also go by the name \emph{partial sums}; the completeness facts behind general convergence arguments belong to a later chapter of calculus, and none of them are needed for the simple limit here.)

\begin{envproposition}{Proposition}{A.5}{Infinite geometric series}
If \(\lvert r \rvert < 1\), then for any number \(a\), \[ \sum_{k = 0}^{\infty} a r^{k} = a + ar + ar^{2} + \cdots = \frac{a}{1 - r}, \] and in particular the tail starting at \(k = 1\) sums to \[ \sum_{k = 1}^{\infty} r^{k} = \frac{r}{1 - r}. \]
\end{envproposition}
\begin{envproof}{Proof}{}{}
By Proposition A.4 the \(m\)-th running total is \(a \cdot \dfrac{1 - r^{m+1}}{1 - r}\). As \(m \to \infty\) the power \(r^{m+1}\) tends to \(0\), so the totals tend to \(\dfrac{a}{1 - r}\). The tail formula is the case \(a = r\): \(r + r^{2} + r^{3} + \cdots = \dfrac{r}{1 - r}\). \qedmark
\end{envproof}
\begin{workedexample}
\begin{envexample}{Example}{A.3}{}
Evaluate each infinite sum.

\begin{enumerate}
  \item \(\dfrac{1}{2} + \dfrac{1}{4} + \dfrac{1}{8} + \cdots\)
  \item \(0.999\ldots\), read as \(\dfrac{9}{10} + \dfrac{9}{100} + \dfrac{9}{1000} + \cdots\)
\end{enumerate}
\end{envexample}
\begin{envsolution}{Solution}{}{}
\begin{enumerate}
  \item This is Proposition A.5 with \(a = \tfrac{1}{2}\) and \(r = \tfrac{1}{2}\): \[ \frac{1}{2} + \frac{1}{4} + \frac{1}{8} + \cdots = \frac{1/2}{1 - 1/2} = 1. \] The picture matches the algebra: eat half the cake, then half of what remains, then half of that — the running total approaches one whole cake, and in the limit exactly fills it.
  \item Here \(a = \tfrac{9}{10}\) and \(r = \tfrac{1}{10}\): \[ 0.999\ldots = \frac{9/10}{1 - 1/10} = \frac{9/10}{9/10} = 1. \] Not approximately \(1\) — exactly \(1\). The two decimal expressions \(0.999\ldots\) and \(1\) are two names for the same number, and the infinite-sum reading is what makes that a computation rather than a debate.
\end{enumerate}
\end{envsolution}
\end{workedexample}
\begin{envcaution}{Caution}{}{}
The hypothesis \(\lvert r \rvert < 1\) is doing real work. For \(r = 2\) the series \(1 + 2 + 4 + 8 + \cdots\) grows without bound and has no sum, while the formula, applied blindly, would report \(\frac{1}{1 - 2} = -1\) — nonsense. The formula \(\frac{a}{1 - r}\) does not \emph{create} a sum; it reports the limit of the running totals, and — unless every term is zero — that limit exists only when \(\lvert r \rvert < 1\).
\end{envcaution}
The geometric series is one of two finite-sum patterns worth carrying forward. The other is a short family of \emph{power sums} — the totals of the first \(n\) whole numbers, their squares, and their cubes — which surface the moment Chapter 6 builds the integral as a limit of sums over \(n\) equal pieces. Like the geometric sum, they drop straight out of telescoping.

\begin{envproposition}{Proposition}{A.6}{Power sums}
For every positive integer \(n\), \[ \begin{aligned}
        \sum_{k=1}^{n} k &= \frac{n(n+1)}{2}, \\[2pt]
        \sum_{k=1}^{n} k^{2} &= \frac{n(n+1)(2n+1)}{6}, \\[2pt]
        \sum_{k=1}^{n} k^{3} &= \left[\frac{n(n+1)}{2}\right]^{2}.
      \end{aligned} \]
\end{envproposition}
\begin{envproof}{Proof}{}{}
Each formula comes from the same collapse that gave the geometric sum. Telescoping supplies \(\sum_{k=1}^{n}\bigl[(k+1)^{2} - k^{2}\bigr] = (n+1)^{2} - 1\), since every interior square is added once and subtracted once; and expanding the summand, \((k+1)^{2} - k^{2} = 2k + 1\). Equating the two evaluations, \[ 2\sum_{k=1}^{n} k + n = (n+1)^{2} - 1 = n^{2} + 2n, \] which solves to \(\sum_{k=1}^{n} k = \tfrac{n(n+1)}{2}\). Climbing one power, \(\sum_{k=1}^{n}\bigl[(k+1)^{3} - k^{3}\bigr] = (n+1)^{3} - 1\) with \((k+1)^{3} - k^{3} = 3k^{2} + 3k + 1\) gives \[ 3\sum_{k=1}^{n} k^{2} + 3\sum_{k=1}^{n} k + n = (n+1)^{3} - 1; \] substituting the value of \(\sum k\) just found and solving leaves \(\sum_{k=1}^{n} k^{2} = \tfrac{n(n+1)(2n+1)}{6}\). The cube sum follows by the identical step one power higher: from \((k+1)^{4} - k^{4} = 4k^{3} + 6k^{2} + 4k + 1\), telescoping gives \(4\sum_{k=1}^{n} k^{3} + 6\sum k^{2} + 4\sum k + n = (n+1)^{4} - 1\), and substituting the two sums already found leaves \(4\sum k^{3} = n^{2}(n+1)^{2}\), that is \(\sum_{k=1}^{n} k^{3} = \bigl[\tfrac{n(n+1)}{2}\bigr]^{2}\). \qedmark
\end{envproof}
A quick check pins the shapes down: at \(n = 3\) the square sum reads \(1 + 4 + 9 = 14\), and \(\tfrac{3 \cdot 4 \cdot 7}{6} = 14\) agrees. What matters is less any single value than the whole shape at once — to add the areas of \(n\) thin rectangles under \(y = x^{2}\) and watch the total settle as \(n \to \infty\), one needs \(\sum k^{2}\) as a formula in \(n\), not a number. That is exactly the step behind the first integral computed from its definition in Chapter 6.

Two formulas from this section will see the most reuse: the geometric tail bound \(\sum_{k \ge 1} r^{k} = \frac{r}{1 - r}\) of Proposition A.5, behind any argument that traps a series under a geometric tail; and the square sum \(\sum_{k=1}^{n} k^{2} = \frac{n(n+1)(2n+1)}{6}\) of Proposition A.6, behind the first integrals Chapter 6 computes from scratch.


\sechead{A.4}{Partial Fractions}
Adding two fractions over a common denominator is routine; the reverse — breaking one complicated fraction back into simple pieces — is \emph{partial fractions}. A rational function is a quotient of polynomials \(P(x)/Q(x)\), and partial fractions rewrites it as a sum of fragments simple enough to handle one at a time. The smallest instance is worth seeing first: \[ \frac{1}{x(x + 1)} = \frac{1}{x} - \frac{1}{x + 1}, \] as recombining the right-hand side over a common denominator confirms. (Read left to right, this is exactly the split that makes \(\sum \tfrac{1}{k(k + 1)}\) telescope — §A.3’s collapsing sum in disguise.) Its real payoff is in integration, where each fragment turns out to have a standard antiderivative.

\begin{envproposition}{Proposition}{A.7}{Partial-fraction decomposition}
Let \(P(x)/Q(x)\) be a \emph{proper} rational function — one whose numerator has lower degree than its denominator — and factor \(Q\) completely into linear and irreducible quadratic factors. Then \(P/Q\) equals a sum of partial fractions, contributed factor by factor:

\begin{itemize}
  \item each distinct linear factor \((x - a)\) contributes one term \(\dfrac{A}{x - a}\);
  \item each repeated linear factor \((x - a)^{m}\) contributes \(\dfrac{A_{1}}{x - a} + \dfrac{A_{2}}{(x - a)^{2}} + \cdots + \dfrac{A_{m}}{(x - a)^{m}}\);
  \item each irreducible quadratic factor \((x^{2} + px + q)\) — one with no real root, i.e. \(p^{2} < 4q\), so it cannot be split into real linear factors — contributes \(\dfrac{Bx + C}{x^{2} + px + q}\); a repeated one \((x^{2} + px + q)^{m}\) contributes one such term over each power up to \(m\).
\end{itemize}
In every case the fragment’s numerator sits one degree below its denominator — a constant over a linear factor, a linear \(Bx + C\) over a quadratic — and a factor raised to a power gets one fragment for each power up to that exponent, exactly the freedom needed to match an arbitrary proper numerator. The constants are then uniquely determined by \(P\) and \(Q\).
\end{envproposition}
That such a decomposition always exists, and with only one choice of constants, is a fact of linear algebra: matching numerators forces a system of linear equations in the unknowns that has exactly one solution. We take the \emph{shape} as given and spend the effort where the work really is — pinning down the constants.

\begin{envstrategy}{Strategy}{A.2}{Finding the constants}
\begin{steps}
  \item If the fraction is \emph{improper} (numerator degree \(\ge\) denominator degree), divide first: \(P/Q = (\text{a polynomial}) + (\text{proper remainder})/Q\), and decompose only the remainder.
  \item Factor the denominator completely and write the template from Proposition A.7 with the constants left as unknowns.
  \item Multiply through by \(Q(x)\) to clear denominators, then solve for the constants — either by matching coefficients of like powers of \(x\), or by substituting convenient values of \(x\) (a factor’s own root kills every other term at once).
\end{steps}
\begin{informal}
The substitution trick in step 3 is the \emph{cover-up} shortcut: to get the constant over a \emph{distinct} linear factor \((x - a)\), set \(x = a\) everywhere except in that one factor. It reaches only the top constant of a repeated factor, and an irreducible quadratic has no real root to substitute at all — those remaining constants come from matching coefficients.
\end{informal}
\end{envstrategy}
\begin{workedexample}
\begin{envexample}{Example}{A.4}{}
Decompose \(\dfrac{5x - 1}{(x - 1)(x + 2)}\).
\end{envexample}
\begin{envsolution}{Solution}{}{}
Two distinct linear factors call for the template \(\dfrac{5x - 1}{(x - 1)(x + 2)} = \dfrac{A}{x - 1} + \dfrac{B}{x + 2}\). Multiplying through by \((x - 1)(x + 2)\) gives \[ 5x - 1 = A(x + 2) + B(x - 1). \] Setting \(x = 1\) leaves \(4 = 3A\), so \(A = \tfrac{4}{3}\); setting \(x = -2\) leaves \(-11 = -3B\), so \(B = \tfrac{11}{3}\). Hence \[ \frac{5x - 1}{(x - 1)(x + 2)} = \frac{4/3}{x - 1} + \frac{11/3}{x + 2}. \] Each substitution was a root of one factor, which zeroed the other term and handed back a single constant — the cover-up shortcut in action.
\end{envsolution}
\end{workedexample}
\begin{workedexample}
\begin{envexample}{Example}{A.5}{}
Decompose each.

\begin{enumerate}
  \item \(\dfrac{x}{(x - 1)^{2}}\), with a repeated linear factor.
  \item \(\dfrac{x + 1}{x(x^{2} + 1)}\), with an irreducible quadratic factor.
\end{enumerate}
\end{envexample}
\begin{envsolution}{Solution}{}{}
\begin{enumerate}
  \item The repeated factor \((x - 1)^{2}\) needs both powers: \(\dfrac{x}{(x - 1)^{2}} = \dfrac{A}{x - 1} + \dfrac{B}{(x - 1)^{2}}\). Clearing denominators, \(x = A(x - 1) + B\). Setting \(x = 1\) gives \(B = 1\); matching the coefficient of \(x\) gives \(A = 1\). So \(\dfrac{x}{(x - 1)^{2}} = \dfrac{1}{x - 1} + \dfrac{1}{(x - 1)^{2}}\).
  \item The quadratic \(x^{2} + 1\) does not factor over the reals — there is no real root to cover up — so it carries a linear numerator, and its constants come from matching coefficients: \(\dfrac{x + 1}{x(x^{2} + 1)} = \dfrac{A}{x} + \dfrac{Bx + C}{x^{2} + 1}\). Clearing denominators, \(x + 1 = A(x^{2} + 1) + (Bx + C)x = (A + B)x^{2} + Cx + A\). Matching powers: the constant term gives \(A = 1\), the coefficient of \(x\) gives \(C = 1\), and the coefficient of \(x^{2}\) gives \(A + B = 0\), so \(B = -1\). Hence \(\dfrac{x + 1}{x(x^{2} + 1)} = \dfrac{1}{x} + \dfrac{1 - x}{x^{2} + 1}\).
\end{enumerate}
\end{envsolution}
\end{workedexample}
\begin{envcaution}{Caution}{}{}
Partial fractions applies only to a \emph{proper} fraction. Faced with \(\dfrac{x^{2}}{x - 1}\), divide first: \(x^{2} = (x - 1)(x + 1) + 1\), so \(\dfrac{x^{2}}{x - 1} = x + 1 + \dfrac{1}{x - 1}\). No combination of the form \(\dfrac{A}{x - 1}\) alone can reproduce a numerator that outgrows its denominator — the polynomial part must be divided out before the templates apply.
\end{envcaution}
Every fragment this method produces — \(\dfrac{A}{x - a}\), \(\dfrac{A}{(x - a)^{k}}\), and \(\dfrac{Bx + C}{(x^{2} + px + q)^{k}}\) — has a standard antiderivative, which is why partial fractions is the standard opening move for integrating a rational function. The algebra that gets you there is complete here.


\sechead{A.5}{The Conic Sections}
Slice a cone with a flat plane and the rim of the cut is a \emph{conic section}: a circle, an ellipse, a parabola, or a hyperbola, according to the angle of the slice. Each also has a cleaner description in terms of distance, and each has a \emph{standard equation} — the simplest form, obtained by placing the curve’s centre or vertex at the origin. This tour records the four standard forms and the features to read straight off them.

\runin{The circle.} A circle is the set of points a fixed distance \(r\) — the \emph{radius} — from a centre; it is the conic already familiar from precalculus, and the template the other three vary. Centred at the origin, the condition \(\sqrt{x^2 + y^2} = r\) squares to the standard equation \[ x^2 + y^2 = r^2. \]

\begin{figureblock}
  \includegraphics[width=107.78mm]{appA/figs/conic-circle}
\figcaption{Figure A.1}{A circle of radius \(r\) centred at the origin: \(x^2 + y^2 = r^2\).}
\end{figureblock}
\runin{The parabola.} Fix a point \(F\) (the \emph{focus}) and a line (the \emph{directrix}). A parabola is the set of points equidistant from the two. Take focus \(F = (0, p)\) with \(p > 0\) and directrix \(y = -p\). Equating a point’s distance to the focus with its distance to the directrix — the latter being the vertical gap \(y + p\) down to the line \(y = -p\) — gives \(x^2 + (y - p)^2 = (y + p)^2\), which simplifies to the standard equation \[ x^2 = 4py, \] a curve opening upward from its vertex at the origin. (Putting the focus on the \(x\)-axis instead gives the sideways form \(y^2 = 4px\).)

\begin{figureblock}
  \includegraphics[width=150.53mm]{appA/figs/conic-parabola}
\figcaption{Figure A.2}{A parabola \(x^2 = 4py\), with focus \(F = (0, p)\) and directrix \(y = -p\).}
\end{figureblock}
\runin{The ellipse.} Fix two points \(F_1, F_2\) (the \emph{foci}). An ellipse is the set of points whose distances to the two foci add to a fixed total. Placing the foci at \((\pm c, 0)\) and writing the total as \(2a\) leads to the standard equation \[ \frac{x^2}{a^2} + \frac{y^2}{b^2} = 1, \qquad b^2 = a^2 - c^2, \] with \(a > b > 0\). The number \(b\) has a clean meaning: at the top of the ellipse, \((0, b)\), symmetry makes the two focal distances equal, and since they sum to \(2a\) each equals \(a\); the right triangle with legs \(b\) and \(c\) and hypotenuse \(a\) then gives \(a^2 = b^2 + c^2\) — the relation above. The curve crosses the axes at \((\pm a, 0)\) and \((0, \pm b)\), and the longer semi-axis \(a\) runs along the line through the foci. Because the foci sit \emph{inside} the curve, \(c < a\), and \(c = \sqrt{a^2 - b^2}\) places them between the centre and the vertices.

\begin{figureblock}
  \includegraphics[width=150.53mm]{appA/figs/conic-ellipse}
\figcaption{Figure A.3}{An ellipse \(\frac{x^2}{a^2} + \frac{y^2}{b^2} = 1\) with foci \(F_1, F_2\) at \((\pm c, 0)\), where \(c = \sqrt{a^2 - b^2}\).}
\end{figureblock}
\runin{The hyperbola.} Keep two foci, but now ask for a fixed \emph{difference} of distances rather than a sum. With foci at \((\pm c, 0)\) and the difference written \(2a\), the standard equation is \[ \frac{x^2}{a^2} - \frac{y^2}{b^2} = 1, \qquad b^2 = c^2 - a^2. \] The curve has two branches with vertices at \((\pm a, 0)\), and far from the centre it draws arbitrarily close to the two \emph{asymptotes} \(y = \pm\frac{b}{a}x\) — the crossed lines whose slopes \(\pm b/a\) set how widely the branches open. This time the foci lie \emph{beyond} the vertices, so \(c > a\), and the relation turns into \(c = \sqrt{a^2 + b^2}\): a \emph{plus} where the ellipse had a minus, precisely because the focus has moved from inside the curve to outside it.

\begin{figureblock}
  \includegraphics[width=150.53mm]{appA/figs/conic-hyperbola}
\figcaption{Figure A.4}{A hyperbola \(\frac{x^2}{a^2} - \frac{y^2}{b^2} = 1\), with vertices \((\pm a, 0)\), foci \((\pm c, 0)\), and asymptotes \(y = \pm\frac{b}{a}x\).}
\end{figureblock}
\begin{workedexample}
\begin{envexample}{Example}{A.6}{}
Identify the curve \(4x^2 + 9y^2 = 36\) and locate its foci.
\end{envexample}
\begin{envsolution}{Solution}{}{}
Divide through by \(36\) to reach standard form: \[ \frac{x^2}{9} + \frac{y^2}{4} = 1. \] Both squared terms are added, so the curve is an ellipse, with \(a^2 = 9\) and \(b^2 = 4\) — that is \(a = 3\) and \(b = 2\). Since \(a > b\), the major axis is horizontal, and \(c = \sqrt{a^2 - b^2} = \sqrt{9 - 4} = \sqrt{5}\). The foci therefore sit at \((\pm\sqrt{5},\, 0)\), about \((\pm 2.24,\, 0)\). Had the sign between the terms been a minus, the same reading would have named a hyperbola instead, with \(c = \sqrt{a^2 + b^2}\).
\end{envsolution}
\end{workedexample}
Those are the four curves, their standard equations, and the features — centre, radius, focus, directrix, foci, axes, and asymptotes — that each equation puts on display. Moving a centre or vertex off the origin only replaces \(x\) and \(y\) with \((x - h)\) and \((y - k)\); the shapes, and everything this tour named, are unchanged.

Read backwards, that last remark is a tool. A second-degree equation handed over in \emph{expanded} form — \(x^2 + y^2 - 6x + 4y = 12\), say — hides its centre until you \emph{complete the square} in each variable, grouping each pair and adding the square of half the linear coefficient to both sides: \[ \begin{aligned}
    &(x^2 - 6x + 9) + (y^2 + 4y + 4) = 12 + 9 + 4, \\
    &\text{that is}\qquad (x - 3)^2 + (y + 2)^2 = 25,
  \end{aligned} \] a circle of radius \(5\) centred at \((3, -2)\). The same single move, \(x^2 + px + q = \left(x + \tfrac{p}{2}\right)^2 + \left(q - \tfrac{p^2}{4}\right)\), is what later lets an integral such as \(\int \dfrac{dx}{x^2 + px + q}\) be recentred and evaluated (Chapter 8). Completing the square is the one algebraic step that turns such an expression — one with no \(xy\) cross-term — into a standard one; a genuine \(xy\) term would call for a rotation of the axes as well, a step beyond this tour.


\sechead{A.6}{Absolute Value and Inequalities}
The absolute value \(|x|\) is the size of \(x\) with its sign discarded — \(|3| = 3\) and \(|-3| = 3\). Its real use in calculus is geometric: \(|x - a|\) is the \emph{distance} between \(x\) and \(a\) on the number line, and every \(\varepsilon\)–\(\delta\) statement is a sentence about such distances. Reading absolute-value inequalities fluently, in both directions, is therefore a standing prerequisite rather than a one-off trick.

The single most useful reading turns a distance bound into an interval. For \(\delta > 0\), \[ |x - a| < \delta \quad\Longleftrightarrow\quad a - \delta < x < a + \delta, \] since “within \(\delta\) of \(a\)” is exactly “no further than \(\delta\) to either side.” The set on the right is the \emph{\(\delta\)-neighbourhood} of \(a\): an open interval of width \(2\delta\) centred at \(a\). Closing the inequality closes the interval — \(|x - a| \le \delta \Leftrightarrow a - \delta \le x \le a + \delta\) — and reversing it turns the region inside out: \(|x - a| > \delta\) describes the two rays \emph{outside}, \(x < a - \delta\) or \(x > a + \delta\). This is the unpacking behind every limit and continuity statement, and behind the interval of convergence of a power series, where \(|x - a| < R\) reads as “\(x\) lies within \(R\) of the centre \(a\).”

\begin{workedexample}
\begin{envexample}{Example}{A.7}{}
Solve \(|2x - 3| < 5\), and describe the solution as a neighbourhood.
\end{envexample}
\begin{envsolution}{Solution}{}{}
Unpack the bound into a double inequality and work the two sides together: \[ \begin{aligned}
          |2x - 3| < 5
            &\;\Longleftrightarrow\; -5 < 2x - 3 < 5 \\
            &\;\Longleftrightarrow\; -2 < 2x < 8 \\
            &\;\Longleftrightarrow\; -1 < x < 4.
        \end{aligned} \] The solution is the open interval \((-1, 4)\). Its midpoint is \(\tfrac{-1 + 4}{2} = \tfrac{3}{2}\) and its half-width is \(\tfrac{5}{2}\), so the same set reads \(\left|x - \tfrac{3}{2}\right| < \tfrac{5}{2}\) — a neighbourhood of \(\tfrac{3}{2}\). Dividing the original bound by \(2\) shows why: \(|2x - 3| = 2\left|x - \tfrac{3}{2}\right|\), so \(|2x - 3| < 5\) is the same demand as \(\left|x - \tfrac{3}{2}\right| < \tfrac{5}{2}\).
\end{envsolution}
\end{workedexample}
One inequality outranks all others for the proofs in this book. It bounds the size of a sum by the sizes of its parts, and it is the lever behind almost every \(\varepsilon\)–\(\delta\) estimate.

\begin{envproposition}{Proposition}{A.8}{Triangle inequality}
For all real numbers \(x\) and \(y\), \[ |x + y| \le |x| + |y|, \] and consequently \(\bigl||x| - |y|\bigr| \le |x - y|\).
\end{envproposition}
\begin{envproof}{Proof}{}{}
Every real number lies between the negative and the positive of its own size: \(-|x| \le x \le |x|\), and likewise \(-|y| \le y \le |y|\). Adding the two chains term by term, \[ -\bigl(|x| + |y|\bigr) \le x + y \le |x| + |y|, \] which says exactly that \(x + y\) is no further than \(|x| + |y|\) from zero — that is, \(|x + y| \le |x| + |y|\). For the second form, apply the first to \(|x| = |(x - y) + y| \le |x - y| + |y|\), so that \(|x| - |y| \le |x - y|\); swapping \(x\) and \(y\) gives \(|y| - |x| \le |x - y|\) as well, and the two together are precisely \(\bigl||x| - |y|\bigr| \le |x - y|\). \qedmark
\end{envproof}
In practice the triangle inequality appears whenever an error must be split and controlled piece by piece: to force \(|f(x) - L|\) below \(\varepsilon\), bound it by a sum of quantities each held under \(\tfrac{\varepsilon}{2}\), and the total stays under \(\varepsilon\). The chapters that prove limits, continuity, and convergence lean on this move constantly.

\begin{envcaution}{Caution}{}{}
An absolute-value inequality is \emph{not} disposed of by simply dropping the bars. The statement \(|x - a| < \delta\) unpacks to the \emph{pair} \(-\delta < x - a < \delta\), not to \(x - a < \delta\) alone — the lower bound is half of what was said. Keeping only the upper half retains the points on one side of \(a\) and silently discards the rest of the neighbourhood.
\end{envcaution}\end{document}
